% Study Characteristics Subsection
\subsection{Study Characteristics}

The 13 studies included in this review exhibit substantial variation in sample size, study design, and validation setting. Detection studies enrolled a median of 44 participants (range 1 to 384), while forecasting studies enrolled a median of 11 participants (range 6 to 70). The total sample comprises 687 participants in detection studies and 225 in forecasting studies.

\subsubsection{Study Design}

Only two studies employed prospective designs. \textcite{spahrDeepLearningbasedDetection2025} conducted a multi-center prospective study across eight epilepsy monitoring units with 384 patients. \textcite{dongTwoLayerEnsembleMethod2022} prospectively monitored 68 patients at home for up to three months. The remaining studies used retrospective analysis of existing datasets or small-scale feasibility designs.

Three detection studies used retrospective hospital data. \textcite{reintjesECGBasedDetectionEpileptic2025} analyzed the SeizeIT2 dataset with 120 patients and 856 seizures from real-world wearable ECG recordings. \textcite{wangEpilepticSeizureDetection2025} analyzed 28 patients with 62 seizures from hospital monitoring. \textcite{elemamAutomatedValidatedTool2025} conducted a cross-sectional study with 198 patients completing a questionnaire and 30 patients providing HRV data.

Two studies represent early-phase validation. \textcite{fineDetectionKeyAutomated2025} reported a Phase 1 study with 15 patients in training and 3 patients in independent testing. \textcite{singhrathoreDevelopmentMultimodalMachine2024} presented a single-patient case study with 6 hours of recording. \textcite{borujenyDetectionEpilepticSeizure2013}, the earliest study in this review, evaluated only 3 patients with 20 seizures.

\subsubsection{Forecasting Study Designs}

All five forecasting studies employed retrospective designs. \textcite{vielufSeizureMonitoringCombined2025} analyzed data from a Phase IV clinical trial with 70 patients and 5437 patient-days. \textcite{meiselMachineLearningWristband2020} used leave-one-subject-out cross-validation on 69 patients with 452 seizures. \textcite{stirlingForecastingSeizureLikelihood2021} conducted retrospective analysis with pseudo-prospective evaluation on 11 patients tracked for a mean of 14.6 months. \textcite{nasseriAmbulatorySeizureForecasting2021} analyzed 6 patients with temporal split validation. \textcite{odeDevelopmentEpilepticSeizure2023} studied 66 patients with 85 seizures using patient-specific anomaly detection.

\subsubsection{Temporal Trends}

Research activity has accelerated over time. Only one study was published between 2013 and 2015. Three forecasting studies were published between 2020 and 2022. Nine detection studies were published between 2023 and 2026. This pattern suggests increased research attention on detection approaches in recent years.

\subsubsection{Setting Distribution}

Eight studies were conducted in hospital or EMU settings. Six studies included home or ambulatory validation. \textcite{dongTwoLayerEnsembleMethod2022} achieved the longest home monitoring duration with 788 overnight recordings over three months. \textcite{stirlingForecastingSeizureLikelihood2021} achieved the longest overall monitoring period with a mean of 14.6 months per patient using consumer smartwatches.

\textbf{Key limitation.} The evidence base includes four studies with sample sizes below 20 participants. \textcite{singhrathoreDevelopmentMultimodalMachine2024} is a single-patient case study with limited generalizability. \textcite{borujenyDetectionEpilepticSeizure2013} evaluated only 3 patients. Small sample sizes limit the reliability of performance estimates and generalizability to broader populations.
